\chapter{A beam coupled to an oscillator}\label{ch:beam}

This case study couples a partial differential equation (a beam, with
infinitely many modes) to an ordinary one (a mass-spring-damper), derives an
entire characteristic function and transfer functions, and computes the
poles without modal truncation. The results are validated with an
independent finite-element model.

\begin{figure}[H]
\centering
\begin{tikzpicture}[>=Latex,thick,font=\small,
  spring/.style={decorate,decoration={zigzag,segment length=3.6pt,amplitude=3pt}}]
  \foreach \offset/\labeltext in {0/{(a) Absorber: independent coordinates},8/{(b) Tip mass, grounded spring and damper}} {
    \begin{scope}[xshift=\offset cm]
      \node[anchor=west,font=\small\bfseries] at (-.3,1.1) {\labeltext};
      \fill[pattern=north east lines] (-.4,-.45) rectangle (0,.45);
      \draw[very thick] (0,-.45)--(0,.45);
      \fill[blue!15] (0,-.09) rectangle (4.4,.09);
      \draw (0,-.09)--(4.4,-.09) (0,.09)--(4.4,.09);
      \node at (2.1,.43) {$EI,\ \rho A,\ L$};
      \draw[->,thin] (.25,-.55)--(1.8,-.55) node[right] {$x$};
      \draw[->,blue!70!black] (2.5,-.1)--(2.5,-1.05) node[below] {$w(x,t)$};
    \end{scope}
  }
  % Independent oscillator linked to a massless tip crossbar.
  \fill (4.4,0) circle (2pt);
  \draw (4.4,0)--(4.4,-.4) (3.75,-.4)--(5.05,-.4);
  \draw[spring] (3.75,-.4)--(3.75,-1.75);
  \node[left] at (3.5,-1.05) {$k$};
  \draw (5.05,-.4)--(5.05,-.8);
  \draw (4.81,-1.25)--(4.81,-.8)--(5.29,-.8)--(5.29,-1.25);
  \draw (4.84,-1.02)--(5.26,-1.02) (5.05,-1.02)--(5.05,-1.75);
  \node[right] at (5.35,-1.05) {$c$};
  \draw (3.75,-1.75)--(5.05,-1.75) (4.4,-1.75)--(4.4,-1.98);
  \draw[fill=orange!20] (3.65,-1.98) rectangle (5.15,-2.65);
  \node at (4.4,-2.3) {$m,\ z(t)$};
  \draw[->,blue!70!black] (5.8,.1)--(5.8,-.5) node[right] {$y=w(L,t)$};
  \draw[->,orange!80!black] (5.65,-2.0)--(5.65,-2.85) node[right] {$z$};
  \draw[->,red!70!black] (4.4,.85)--(4.4,.15) node[midway,right] {$f$};
  \node[align=center,font=\footnotesize] at (3,-3.3)
    {Spring and damper act on $z-y$.\\The mass is not connected to the ground.};
  % Rigidly attached point mass, spring and damper connected to ground.
  \begin{scope}[xshift=8cm]
    \draw[fill=orange!20] (3.8,-.3) rectangle (5,.3);
    \node at (4.4,0) {$m$};
    \draw (4.4,-.3)--(4.4,-.55) (3.75,-.55)--(5.05,-.55);
    \draw[spring] (3.75,-.55)--(3.75,-1.9);
    \node[left] at (3.5,-1.2) {$k$};
    \draw (5.05,-.55)--(5.05,-.95);
    \draw (4.81,-1.4)--(4.81,-.95)--(5.29,-.95)--(5.29,-1.4);
    \draw (4.84,-1.17)--(5.26,-1.17) (5.05,-1.17)--(5.05,-1.9);
    \node[right] at (5.35,-1.2) {$c$};
    \draw (3.2,-1.9)--(5.6,-1.9);
    \fill[pattern=north east lines] (3.2,-2.13) rectangle (5.6,-1.9);
    \draw[->,blue!70!black] (5.8,.1)--(5.8,-.65) node[right] {$y=z$};
    \node[align=center,font=\footnotesize] at (3,-3.3)
      {Spring and damper act on $y$.\\The mass moves rigidly with the tip.};
  \end{scope}
\end{tikzpicture}

\caption{The two topologies. The connection transmits transverse force but
no moment; displacements are positive downwards.}
\label{fig:schematic}
\end{figure}

\section{Model}

A uniform Euler--Bernoulli beam of length $L$, bending stiffness $EI$ and
mass per length $\mu_b=\rho A$ is clamped at $x=0$. In topology (a), a mass
$m$ with displacement $z(t)$ is connected to the tip displacement
$y(t)=w(L,t)$ by a spring $k$ and a damper $c$ in parallel; the mass is not
grounded (a tuned-mass absorber). The beam is elastic; all dissipation is in
$c$. Shear deformation, rotary inertia and gravity (already balanced) are
neglected:
\begin{gather}
  \mu_bw_{tt}+EIw_{xxxx}=0,\quad 0<x<L,\qquad w(0,t)=w_x(0,t)=0,\\
  EIw_{xx}(L,t)=0,\qquad -EIw_{xxx}(L,t)=F(t)+f(t),\\
  f=k(z-y)+c(\dot z-\dot y),\qquad m\ddot z+c(\dot z-\dot y)+k(z-y)=u(t),
\end{gather}
where $f$ is the force of the connection on the beam and $F$, $u$ are
external forces at the tip and on the mass.

\section{The bare-beam receptance}

With $w=W(x)e^{st}$ and $\theta^4=q=-\mu_bL^4s^2/EI$, the clamped solution
is $W(x)=a[\cosh bx-\cos bx]+d[\sinh bx-\sin bx]$, $b=\theta/L$. The
zero-moment condition fixes $d/a$, and the shear condition gives the tip
receptance
\[
  H_b(s)=\frac{y}{f}=\frac{L^3}{EI}\frac{A(q)}{B(q)},\quad
  A(q)=\frac{\cosh\theta\sin\theta-\sinh\theta\cos\theta}{\theta^3},\quad
  B(q)=1+\cosh\theta\cos\theta .
\]
$A$ and $B$ are entire in $q$: they are invariant under $\theta\mapsto\ii\theta$.
Near $q=0$, $A(q)=\tfrac23-\tfrac q{315}+\tfrac{q^2}{623700}-\tfrac{q^3}{5108103000}+\cdots$
and $B(0)=2$, so $H_b(0)=L^3/(3EI)$, the static compliance of a cantilever.
Write $H_b=N_b/D_b$ with $N_b=(L^3/EI)A$, $D_b=B$.

\section{Characteristic function and transfer functions}

With $Z(s)=k+cs$ the Laplace-domain equations are
\[
  \begin{bmatrix}D_b+N_bZ & -N_bZ\\ -Z & ms^2+Z\end{bmatrix}
  \begin{bmatrix}y\\ z\end{bmatrix}=
  \begin{bmatrix}N_bF\\ u\end{bmatrix},
\]
whose determinant is the entire characteristic function
\begin{equation}
  \Delta(s)=D_b(s)(ms^2+cs+k)+N_b(s)\,ms^2(cs+k) .
  \label{eq:beamdelta}
\end{equation}
The coupled poles are not the union of the beam and oscillator poles, and
$1+H_bZ_{\mathrm{eff}}$, with $Z_{\mathrm{eff}}=ms^2(cs+k)/(ms^2+cs+k)$, is
not entire; clearing its denominators gives~\eqref{eq:beamdelta}. The
transfer functions are
\[
  \frac yF=\frac{N_b(ms^2+Z)}{\Delta},\qquad
  \frac yu=\frac{N_bZ}{\Delta},\qquad
  \frac zu=\frac{D_b+N_bZ}{\Delta}.
\]
With zero inputs, the energy
$E=\tfrac12\int_0^L(\mu_bw_t^2+EIw_{xx}^2)\,dx+\tfrac12m\dot z^2+\tfrac12k(z-y)^2$
satisfies $\dot E=-c(\dot z-\dot y)^2\le0$, which confirms the sign
conventions and excludes growing modes for positive parameters. In topology
(b) the mass moves with the tip and $\Delta_{\mathrm{tip}}=D_b+N_b(ms^2+cs+k)$;
it is a different system, with potential energy $ky^2/2$.

\section{Implementation}

With $\tau=\sqrt{\mu_bL^4/EI}$, $\lambda=\tau s$, $\mu=m/(\mu_bL)$,
$\kappa=kL^3/EI$ and $\gamma=cL^3/(EI\tau)$, the function evaluated is
$\Delta=B(-\lambda^2)(\mu\lambda^2+\gamma\lambda+\kappa)+A(-\lambda^2)\mu\lambda^2(\gamma\lambda+\kappa)$,
equal to~\eqref{eq:beamdelta} times $L^3/EI$. In normalized units the
complete computation needs only function handles:
\begin{lstlisting}
m = 0.1;  k = 1.2362363368;  c = 0.05;      % L = EI = rho*A = 1
beta = @(s) (-s.^2).^(1/4);
Nb = @beam_numerator;                        % series near 0, see text
Db = @(s) 1 + cosh(beta(s)).*cos(beta(s));
Zc = @(s) k + c*s;
N = @(s) Nb(s).*(m*s.^2 + Zc(s));
D = @(s) Db(s).*(m*s.^2 + Zc(s)) + Nb(s).*m.*s.^2.*Zc(s);
[p, info] = cpoles(ndpair(N, D), [-60 1 -150 150], 'AssumeAnalytic', true);
\end{lstlisting}
The full script is \path{examples/coupled_beam/beam_direct_nd.m}; the
helper \path{examples/models/coupled_beam_model.m} builds the dimensional
model and the three transfer functions.

\section{Reference case and parameter studies}

The reference case uses $m=0.1$, $k=m\,\omega_{b1}^2=\BaselineStiffness$ with
the first bare-beam frequency $\omega_{b1}=\FirstBareFrequency$, and $c=0.05$.
The ten poles in $[-60,1]\times[-150,150]$ are (upper half-plane):
\begin{center}\small\BaselinePoleTable\end{center}
For $c=0$ the two lowest frequencies are $2.5703$ and $4.7801$: the coupling
splits the original resonance; in the first mode the mass follows the tip,
in the second it moves against it (Figure~\ref{fig:modes}).

\begin{figure}[H]
\centering
\includegraphics[width=\textwidth]{coupled_modes.png}
\caption{Conservative coupled modes from the exact spatial solution,
normalized by the tip displacement.}
\label{fig:modes}
\end{figure}

One parameter at a time was varied: $m\in[0.02,0.5]$ (25 values),
$k\in[0.05,5]$ (36 values) and $c\in[0,1.5]$ (43 values). Every one of the
104 cases has ten poles in the window, each found with a fresh contour
count (previous poles were only seeds).

\begin{figure}[H]
\centering
\includegraphics[width=\textwidth]{pole_sweeps.png}
\caption{Low-frequency poles as the mass, the stiffness and the damping vary.}
\label{fig:sweeps}
\end{figure}

\begin{center}\small\DampingPoleTable\end{center}

Increasing $c$ can make the second pair real. Among the sampled values the
rightmost pole is leftmost at $c=\BestDamping$ (abscissa $\BestAlpha$); at
$c=1.5$ the dominant pair returns to $-0.0840\pm2.9549\ii$. Damping in the
connection restrains the very relative motion it dissipates, so more
damping is not always better. This is a statement about the sampled values
and the stated window, not a global optimum.

\begin{figure}[H]
\centering
\includegraphics[width=\textwidth]{trends_and_validation.png}
\caption{Lowest branches versus mass and stiffness, window abscissa versus
damping, and convergence of the finite-element model.}
\label{fig:trends}
\end{figure}

\section{Independent validation}\label{sec:beamvalidation}

A cubic-Hermite finite-element model with consistent mass matrix and an
extra degree of freedom for the absorber mass (\code{coupled\_beam\_fem.m})
is solved as a quadratic eigenvalue problem, without using any result of the
nonrational solver. The largest error among the ten reference poles
decreases with the mesh:
\begin{center}\small\FEMTable\end{center}
Further checks: the static compliance $H_b(0)=1/3$; invariance of the
dimensionless characteristic function under dimensional rescaling; the
bare-beam limit; and the decoupled limit $k=c=0$, where the characteristic
function has a double root at $0$ (the free mass) that is \emph{not} a pole
of the tip transfer function $y/F$ --- a hidden mode. Enlarging the window
to $[-90,1]\times[-300,300]$ at three damping values leaves the rightmost
pole unchanged. The beam is elastic and has no finite accumulation point,
unlike the Kelvin--Voigt beam of Section~\ref{sec:kelvinvoigt}.
